\section{Places of small discrepancy}
\label{sec:low-places}

Fix a model $T$ in the sequence of Theorem~\ref{thm:termination}.
Proposition~\ref{prop:comparison} gives log discrepancy at least one
for every exceptional place over $T$.  We will show that, apart from a
finite set of places, an exceptional place of discrepancy below two is
centred on a surface in a positive boundary component.  Its discrepancy
is then bounded below by two minus that component's coefficient.
This description will restrict the possible divisors used to count
flipped surfaces.

\subsection{The nef trace and local estimates}

Write $B_T$ and $M_T$ for the boundary and the trace of $\bM$ on $T$.
They are $\Q$-Cartier: each is a finite rational combination of globally
defined prime Weil divisors, to which global Weil $\Q$-factoriality applies.
Choose a smooth model
\[
 r\colon W\longrightarrow T
\]
that is projective over $T$ and carries the fixed nef data, as supplied by
Proposition~\ref{prop:common-models}.  Define the actual divisors
\begin{equation}
 \Delta_W=r^*(K_T+B_T+M_T)-K_W-M_W,
 \qquad J_W=r^*M_T-M_W.
 \label{eq:boundary-and-defect}
\end{equation}
Thus the coefficient of a prime $E$ in $\Delta_W$ is
$1-a(E;T,B_T+\bM)$.

\begin{lemma}[Effectivity of the nef-trace defect]
\label{lem:nef-defect}
For every smooth model $r\colon W\to T$ that is projective over $T$ and
carries the fixed nef data, the divisor $J_W$ in
\eqref{eq:boundary-and-defect} is an effective
$r$-exceptional $\Q$-Cartier divisor.
\end{lemma}

\begin{proof}
The pushforward of both $r^*M_T$ and $M_W$ is $M_T$, so $J_W$ is
$r$-exceptional.  For every curve $C$ contracted by $r$,
\[
 (-J_W)\cdot C=M_W\cdot C\geq 0.
\]
To see the last inequality, let $s\colon W\to X'$ be the morphism to the
fixed carrier.  Since $M_W=s^*M'$, its degree on $C$ is zero if $s(C)$ is a
point, and otherwise is the degree of $M'$ on $s(C)$ multiplied by the
degree of the induced map of normalized curves.  Analytic nefness of $M'$
implies that these curve degrees are nonnegative: restrict a sequence of
K\"ahler classes converging to $c_1(M')$ to the normalized curve and pass to
the limit.  We use this consequence of analytic nefness only for the
relative degree calculation above.  Lemma~\ref{lem:negativity}, applied to
$J_W$ and $r_*J_W=0$, now gives $J_W\geq0$.
\end{proof}

We will use the following local estimate both to check ordinary terminality
and to locate places of small generalized discrepancy.  It applies to local
analytic places as well as to places represented by divisors on compact
models.

\begin{lemma}[A local differential estimate]
\label{lem:jacobian}
Let $W$ be a smooth complex space of dimension $n$, and let $E$ be a prime
divisor on a smooth proper birational model of an open subset of $W$.
Suppose its centre $C$ has codimension $c$.  At a general smooth point of
$C$, choose coordinates $x_1,\ldots,x_n$ with
$C=(x_1=\cdots=x_c=0)$.  Writing $v=\operatorname{ord}_E$, one has
\begin{equation}
 a(E;W,0)\geq\sum_{j=1}^{c}v(x_j)\geq c.
 \label{eq:jacobian-bound}
\end{equation}
If $C$ is contained in exactly one positive component $H$ of a rational
divisor $\Delta$, and $H$ is smooth at the general point of $C$ with
coefficient $d\in(0,1)$, choose the coordinates so that $H=(x_1=0)$.
Then
\begin{equation}
 a(E;W,\Delta)
 \geq (1-d)v(x_1)+\sum_{j=2}^{c}v(x_j)
 \geq c-d.
 \label{eq:one-positive-bound}
\end{equation}
If $C$ is contained in no positive component of $\Delta$, then
$a(E;W,\Delta)\geq c$.
\end{lemma}

\begin{proof}
Compute at a general smooth point of $E$ above the chosen coordinate
neighbourhood.  If $t$ is a local equation of $E$, write
$x_j=t^{v(x_j)}u_j$ for $1\leq j\leq c$, where $u_j$ is a unit at the
general point of $E$.  Each $v(x_j)$ is a positive integer.  In expanding
the pullback of $dx_1\wedge\cdots\wedge dx_n$, at most one factor in a
nonzero term can contribute a differential $dt$.  Thus its order along
$E$ is at least $\sum_{j=1}^c v(x_j)-1$.  The remaining coordinate
differentials are holomorphic and do not lower this order.  Adding one to
the coefficient of the relative canonical divisor proves
\eqref{eq:jacobian-bound}.

For the boundary statement, each component with nonpositive coefficient
contributes a nonnegative amount to the log discrepancy.  Components not
containing $C$ have order zero along $E$.  Subtracting just
$d\,v(x_1)$ from \eqref{eq:jacobian-bound} therefore gives
\eqref{eq:one-positive-bound}.  If there is no positive component through
$C$, nothing has to be subtracted.
\end{proof}

\subsection{A finite exceptional set of places}

We next describe the places with generalized log discrepancy below two.
The finite exceptional set is attached to one fixed model; its size need
not be uniform along the sequence.  The distinction according to the
centre of a place follows the corrected discrepancy argument of
Fujino~\cite[Proposition~3.1 and Lemma~3.2]{Fujino2005Addendum}.  We give
the full argument, including the repeated blowups that are needed when a
boundary coefficient exceeds $1/2$.

\begin{lemma}[Places of log discrepancy below two]
\label{lem:low-places}
Let $(T,B_T+\bM)$ be any model of a sequence as in Theorem~\ref{thm:termination}.  There is a
finite set $\mathcal F_T$ of exceptional places over $T$ with the following
property.  If $E$ is exceptional over $T$,
\[
 a(E;T,B_T+\bM)<2,
 \qquad E\notin\mathcal F_T,
\]
then its centre $S$ on $T$ is a surface contained in a unique positive
component $H$ of $B_T$.  The space $T$ and the reduced boundary support are
smooth at the general point of $S$.  If the coefficient of $H$ is $d$,
then
\begin{equation}
 a(E;T,B_T+\bM)\geq2-d.
 \label{eq:low-place-threshold}
\end{equation}
If $B_T=0$, there are only finitely many exceptional places with log
discrepancy below two.
\end{lemma}

\begin{proof}
Choose a smooth modification $r\colon W\to T$, projective over $T$ and
carrying the nef data, such that its exceptional locus is divisorial and, together with
the strict transform of $\Supp B_T$, has simple normal crossings.
Starting with this simple normal crossing divisor, blow up the
intersections of pairs of positive strict transforms until they are
mutually disjoint.  Each intersection is a finite disjoint union of smooth
codimension-two subspaces.  Blowing it up separates the chosen pair,
preserves simple normal crossings, and cannot make a previously separated
pair meet.  Thus finitely many such blowups suffice.

Keep the notation $W$ for the resulting model and $\Delta_W$ for its
boundary in \eqref{eq:boundary-and-defect}.  Its positive components are
exactly the strict transforms of the positive components of $B_T$, with
the same coefficients.  They are smooth and pairwise disjoint.  Every
component exceptional over $T$ has nonpositive coefficient by
generalized canonicity; zero coefficients are allowed.  On a higher
smooth model $u\colon V\to W$,
the nef terms cancel because $M_V=u^*M_W$, giving
\[
 K_V+\Delta_V=u^*(K_W+\Delta_W).
\]
Hence generalized discrepancies over $T$ can be computed as ordinary
discrepancies for $(W,\Delta_W)$, including its negative coefficients.

Put all $r$-exceptional prime divisors on $W$ into $\mathcal F_T$.
There are finitely many, since $W$ is compact.  A further exceptional
place $E$ has centre $C$ on $W$ of codimension at least two: a place with
divisorial centre on a normal model is represented by that divisor.
Lemma~\ref{lem:jacobian} shows that
$a(E;W,\Delta_W)<2$ forces $C$ to have codimension two and to lie in
one of the positive components, say $H$ with coefficient $d$.  The same
lemma gives
\begin{equation}
 a(E;W,\Delta_W)\geq2-d.
 \label{eq:resolution-low-bound}
\end{equation}
If $C$ is not contained in $\Exc(r)$, then $r$ is an isomorphism at its
general point.  Its image is the surface $S$ required in the statement;
smoothness and uniqueness of the positive component descend from $W$.

It remains to treat the centres contained in $\Exc(r)$.  Such a centre
$C$, being codimension two and contained in $H$, must be an irreducible
component of $H\cap F$ for some exceptional divisor $F$ on $W$.
There are only finitely many such components, and each is smooth by
simple normal crossings.  We show that a fixed one contributes only
finitely many places of log discrepancy below two.

Near $C$, the only positive component of $\Delta_W$ is $H$.
Deleting the nonpositive components can only lower discrepancies, so
\[
 a(E;W,dH)\leq a(E;W,\Delta_W)<2
\]
for every place under consideration with centre $C$.  We can therefore
work with the smooth pair $(W,dH)$ and the fixed smooth codimension-two
centre $C\subset H$.

Set $W_0=W$ and $H_0=H$.  Blow up $C$ to obtain $W_1$, with exceptional
divisor $F_1$ and strict transform $H_1$.  If $E$ is not $F_1$, its
centre on $W_1$ is contained in $F_1$ and has codimension at least two.
The crepant boundary on $W_1$ has sole positive component $dH_1$; the
coefficient of $F_1$ is $-(1-d)$.  Lemma~\ref{lem:jacobian} therefore
forces the centre of $E$ to have codimension two and to lie in $H_1$.
It is consequently the entire intersection
\[
 C_1=H_1\cap F_1.
\]
This intersection is a smooth irreducible section of the exceptional
$\mathbb P^1$-bundle over $C$, hence isomorphic to $C$.

Repeat this construction: after $W_j$ has been obtained, blow up
$C_j=H_j\cap F_j$ to obtain $W_{j+1}$, with new exceptional divisor
$F_{j+1}$.  Writing $C_0=C$, the maps and
their centres are
\begin{equation}\label{eq:forced-centres}
 W_j=\operatorname{Bl}_{C_{j-1}}W_{j-1},\qquad
 C_j=H_j\cap F_j\simeq C_{j-1}\qquad(j\geq1).
\end{equation}
The new strict transform
$H_{j+1}$ meets $F_{j+1}$ in a smooth section and is disjoint from the
strict transform of $F_j$.  Inductively, it misses every older exceptional
divisor.  The codimension-two blowup formula gives the coefficient
\begin{equation}
 \operatorname{coeff}_{F_j}(\Delta_j)=-j(1-d),
 \qquad a(F_j;W,dH)=1+j(1-d),
 \label{eq:forced-chain-coefficients}
\end{equation}
where $K_{W_j}+\Delta_j$ is the pullback of $K_W+dH$.
Indeed the new coefficient is the sum of the coefficients of $H_j$ and
$F_j$, minus one.  This is $-(j+1)(1-d)$.

Thus the only way a place of discrepancy below two can remain
unextracted is for its centre to follow these successive intersections.
At every stage, if $E$ has not yet appeared, the argument just given
forces its centre to be $C_j$.  Choose local transverse coordinates
$x,y$ with $H_j=(x=0)$ and $F_j=(y=0)$, and put
$v=\operatorname{ord}_E$.  Retaining the negative coefficient of $F_j$
in the differential estimate gives
\begin{equation}
 \begin{split}
 a(E;W,dH)&\geq (1-d)v(x)+\bigl(1+j(1-d)\bigr)v(y)\\
          &\geq (1-d)+\bigl(1+j(1-d)\bigr)
            =2-d+j(1-d).
 \end{split}
 \label{eq:forced-chain-cutoff}
\end{equation}
For
\[
 J=\left\lceil\frac{d}{1-d}\right\rceil,
\]
the right side is at least two.  Every place with discrepancy below two
and centre $C$ must therefore be one of $F_1,\ldots,F_J$.  Equivalently,
the possible indices in \eqref{eq:forced-chain-coefficients} satisfy
$j<1/(1-d)$; the strict endpoint is included in the cutoff calculation.

All these blowups have global smooth compact analytic centres.  Thus the
divisors just enumerated are global places.  Conversely, an arbitrary
place can be followed on common proper resolutions even when its
original carrying model was not projective; its centres map onto the
preceding centres.  The argument forcing $C_j$ therefore covers every
place over $C$, not only those chosen in advance from this blowup chain.

Add the finitely many divisors obtained for all the finitely many
intersections $H\cap F$ to $\mathcal F_T$.  Every place outside this set
has the centre described earlier, and
\eqref{eq:resolution-low-bound} gives
\eqref{eq:low-place-threshold}.  If the boundary has no positive
component, Lemma~\ref{lem:jacobian} already excludes every further place
of discrepancy below two, proving the last assertion.
The proof used only nonpositivity of the exceptional coefficients on
$W$, so exceptional places of log discrepancy exactly one cause no
additional case.
\end{proof}

\begin{remark}
The second and subsequent blowups in the preceding proof are essential.
For example, when $d=2/3$, the first two divisors have log discrepancies
$4/3$ and $5/3$, and the third has discrepancy two.  This is the phenomenon
in Fujino's correction \cite[Example~2.1]{Fujino2005Addendum}.  The finite
exceptional set in Lemma~\ref{lem:low-places} is defined by centres and
contains the whole finite chain over each exceptional intersection.
\end{remark}
